\begin{claim}{D.42}{Proved}{Order-three pants actions have global multiplier one}
Every element of order three in the full three-cusp
loop-outer group has orientation multiplier
$1\in\widehat\Z^\times$.
Consequently every finite subgroup of that full
group has multiplier image contained in the two
global units $\{1,-1\}$, and its multiplier is
determined by its $2$-adic coordinate.

The closed-surface calculation used in the proof is
also unconditional.  For $g\geq2$ and every prime $p$,
the homomorphism
\[
 \Aut(\widehat\pi_{g,0})
       \longrightarrow\GL(H_1(\widehat\pi_{g,0},\Z_p))
                                                        \tag{D.42a}
\]
is injective on every finite $p$-subgroup.
Here the source consists of actual continuous
automorphisms.  The order-three outer action in the
first assertion is supplied with an actual lift by
the exact rank-two lifting result D.27.
\end{claim}
\begin{proof}
We first prove (D.42a), keeping the two completion
maps and the fixed subgroup explicit.

\emph{The pro-$p$ automorphism calculation.}
Let $P=\pi_{g,0}^{(p)}$ be the pro-$p$ completion.
The closed lower central terms $\gamma_nP$
satisfy
\[
 \bigcap_n\gamma_nP=1,
 \qquad P/\gamma_nP\text{ is torsion-free for every }n.
                                                        \tag{D.42b}
\]
These are precisely
\cite[Propositions 2.14--2.15, pp.~1515--1516]{A-Morales2024},
applied to the orientable surface presentation.
The second proposition checks the primitive degree-two
initial form of its single defining relation; the
first passes the torsion-free nilpotent quotients to
the pro-$p$ completion by the stated canonical quotient
maps.  In particular every
$\gamma_nP/\gamma_{n+1}P$, as a subgroup of
$P/\gamma_{n+1}P$, is torsion-free.

Put
$K_n=\ker(\Aut(P)\to\Aut(P/\gamma_{n+1}P))$ for $n\geq1$.
For $b\in K_n$, the formula
\[
 d_n(b)(\bar z)=b(z)z^{-1}\pmod{\gamma_{n+2}P}
                                                        \tag{D.42c}
\]
defines a continuous homomorphism from $P^{\rm ab}$
to $\gamma_{n+1}P/\gamma_{n+2}P$.
Indeed this last quotient is central modulo
$\gamma_{n+2}P$, which makes the formula additive
in $z$ and trivial on commutators.
An automorphism acting trivially on $P^{\rm ab}$
acts trivially on every lower central graded quotient:
these quotients are generated by the iterated
commutators of its abelianization generators.
Consequently $d_n$ is a group homomorphism on $K_n$,
and its kernel is exactly $K_{n+1}$.
Its target is torsion-free by (D.42b).
A finite-order element of $K_1$ therefore lies
successively in every $K_n$ and is the identity.
Thus the kernel of
$\Aut(P)\to\GL(H_1(P,\Z_p))$ is torsion-free.

\emph{The full fixed subgroup.}
Write $\Pi=\widehat\pi_{g,0}$ and let
$a\in\Aut(\Pi)$ satisfy $a^p=1$ and act trivially
on $H_1(\Pi,\Z_p)$.
Its induced automorphism on $P=\Pi^{(p)}$ is the
identity by the preceding calculation.
There is therefore a specified quotient homomorphism
\[
 E=\Pi\rtimes_a C_p\longrightarrow P\times C_p,
             \qquad(h,s^i)\longmapsto(\bar h,s^i).
                                                        \tag{D.42d}
\]
The natural map $\Pi\to P$ induces an isomorphism
on mod-$p$ cohomology, including products.
Both groups have cohomological $p$-dimension two,
first cohomology $\F_p^{2g}$, and second cohomology
$\F_p$ with the standard nondegenerate surface
cup pairing.  The map is the identity on the
generator characters, so the pairing shows that
it is also an isomorphism in degree two; higher
degrees vanish.  These are the surface duality
comparisons of A.2 and A.7, with the full and
pro-$p$ formations respectively.
Hochschild--Serre comparison for the named map
(D.42d) now gives a ring isomorphism
\[
 H^*(P\times C_p,\F_p)\xrightarrow{\sim}H^*(E,\F_p).
                                                        \tag{D.42e}
\]
In particular $\dim H^n(E,\F_p)=2g+2$ for $n>2$.

Choose a nonzero surface class
$\eta\in H^2(P,\F_p)$ and a nonzero periodic
class $u\in H^2(C_p,\F_p)$.
For every $m\geq1$, the class
$\eta u^m\in H^{2m+2}(E,\F_p)$ is nonzero
under (D.42e).
Every subgroup $S\leq E$ of order $p$ projects
isomorphically to $C_p$, because $\Pi$ is
torsion-free.  Its first coordinate in $P\times C_p$
is trivial because $P$ is torsion-free.
Its normalizer is exactly
\[
               N_E(S)=C_\Pi(S)\times S.                 \tag{D.42f}
\]
An element normalizing $S$ fixes each of its elements
by projection to the abelian quotient $C_p$;
splitting off the uniquely determined $S$ component
proves (D.42f).

Apply the full profinite Brown restriction theorem
\cite[Theorem 2 and \S1.1]{A-KochloukovaZalesskii2007}
with virtual cohomological dimension two.
Its hypotheses hold: $\Pi$ is an open subgroup
of cohomological $p$-dimension two, and there is
no $C_p\times C_p$ in $E$, since every finite
subgroup injects into $E/\Pi=C_p$.
The finite-dimensional source makes its dense
restriction embedding an isomorphism onto a
finite product of the normalizer cohomologies.
Restriction of the nonzero class $\eta u^m$
is nonzero in at least one factor.
For that factor put $C=C_\Pi(S)$.
Because the first coordinate of $S$ in (D.42d)
is trivial, its restriction is exactly
\[
        (\eta|_C)\otimes(u^m|_S)
       \in H^2(C,\F_p)\otimes H^{2m}(S,\F_p).
                                                        \tag{D.42g}
\]
Thus $\eta|_C\ne0$.
The restriction
$H^1(\Pi,\F_p)\to H^1(C,\F_p)$ is injective:
for any nonzero $v$, nondegeneracy supplies $w$
with $v\smile w$ a nonzero multiple of $\eta$,
and its restriction cannot vanish.
Hence $\dim H^1(C,\F_p)\geq2g$ and
$\dim H^2(C,\F_p)\geq1$.
For every normalizer (D.42f), the K\"unneth formula
in degree $n>2$ contributes
\[
  1+\dim H^1(C_\Pi(S),\F_p)
                +\dim H^2(C_\Pi(S),\F_p).              \tag{D.42h}
\]
This holds also for odd $p$: the exterior generator
of $H^*(C_p,\F_p)$ and its degree-two polynomial
generator give one dimension in each nonnegative
degree.  At $p=2$ the degree-one polynomial
generator has the same dimension consequence.
The one factor in (D.42g) already accounts for
all $2g+2$ dimensions.  There is exactly one
conjugacy class of sections, and for its centralizer
\[
 \dim H^1(C,\F_p)=2g,\qquad H^2(C,\F_p)=\F_p.           \tag{D.42i}
\]
Conjugacy under $E$ equals conjugacy under $\Pi$,
since $E=\Pi S$ for each section $S$ and the
$S$ factor centralizes that section.
Consequently these conclusions apply to
$C=C_\Pi(a)$ for the specified section in (D.42d).

We now prove directly that this closed subgroup
equals $\Pi$.
Continuous cohomology has the restriction comparison
\[
 H^j(C,\F_p)=
  \mathop{\rm colim}_{C\leq U\leq_o\Pi}H^j(U,\F_p).
                                                        \tag{D.42j}
\]
For an explicit cochain comparison, choose an
open normal subgroup $V$ of $\Pi$ so small that
the given continuous finite-valued cochain on $C$
factors through $C/(C\cap V)$ in every variable.
The isomorphism $CV/V=C/(C\cap V)$ inflates it to
the open subgroup $CV$.  Cocycle equalities descend
to that finite quotient; for a coboundary choose
a common $V$ also for its primitive cochain.
This proves (D.42j) with the displayed restriction maps.
Each $U$ is the completion of its finite covering
surface, of genus
$g_U=1+[\Pi:U](g-1)$.
For $V\leq U$ the restriction on $H^2$ multiplies
the fundamental class by $[U:V]$.
A nonzero class surviving in (D.42j) in degree two
therefore has a representative at some $U$ for
which every smaller open $V\geq C$ has
$p\nmid[U:V]$.
For these subgroups the restriction on $H^1$
is injective: corestriction composed with restriction
is multiplication by this prime-to-$p$ index.
If $C$ had infinite index, the indices
$[\Pi:V]$ would be unbounded, making
\[
 \dim H^1(V,\F_p)=2+2[\Pi:V](g-1)
\]
unbounded inside the fixed finite-dimensional
$H^1(C,\F_p)$, contrary to (D.42i).
Thus $C$ is open.  The same genus formula and
(D.42i) give $[\Pi:C]=1$.
It follows that $a=1$.
The kernel of (D.42a) on a finite $p$-subgroup
cannot contain an element of order $p$, proving
its asserted injectivity.

\emph{The order-three cyclic cover.}
Let $g$ have order three in the full pants group,
with multiplier $\lambda$.
By D.32 its cusp permutation is a three-cycle.
Label $x,y,z=(xy)^{-1}$ so that it sends their
classes successively to $y,z,x$, with the common
power $\lambda$.
D.27 gives an actual order-three representative
$\alpha\in\Aut(\widehat F(x,y))$.
The equality $\lambda^3=1$ implies $\lambda_3=1$:
an order-three unit of $\Z_3$ reduces to one,
and for $v\in\Z_3^\times$, $r\geq1$, the
valuation of $(1+3^rv)^3-1$ is $r+1$.

Suppose a coordinate $\lambda_q$ were nonidentity.
Then $q\ne3$, and its reduction
$\zeta\in\F_q^\times$ has exact order three;
in particular $q\equiv1\pmod3$ and $q\geq7$.
Define the surjective character
\[
 \chi:F\longrightarrow C_q=\F_q,\qquad
 \chi(x)=\zeta,\quad\chi(y)=1,\quad\chi(z)=\zeta^2.
                                                        \tag{D.42k}
\]
The relation is respected because
$1+\zeta+\zeta^2=0$.
The three scalar identities
$\zeta\cdot1=\zeta$,
$\zeta\cdot\zeta^2=1$, and
$\zeta\cdot\zeta=\zeta^2$
show exactly that $\chi\alpha=\chi$;
all inner conjugators disappear under $\chi$.
Let $N=\ker\chi$.
The specified degree-$q$ cover has one cusp
above each of the original three, since all
three values in (D.42k) are nonzero.
Its type is $(h,3)$ with
$h=(q-1)/2\geq3$, by
$2-2h-3=-q$.
Fill these three cusps and put
\[
 \Pi_h=N/\overline{\langle\!\langle x^q,y^q,z^q
                         \rangle\!\rangle}_N
          =\widehat\pi_{h,0}.                            \tag{D.42l}
\]
Completion of the finite-cover presentation
gives this exact quotient map.
The peripheral powers under $\alpha$ preserve
the indicated closed normal subgroup, so it
induces an actual automorphism $b$ of $\Pi_h$
with $b^3=1$.
Conjugation by $y$ gives the deck automorphism
$T$ of this quotient; $T^q=1$, because $y^q$
has been killed.
The relation $\chi\alpha=\chi$ says that
$bT$ and $Tb$ differ by an inner automorphism.
Their actions on homology therefore commute.

Here is an explicit calculation of that deck
module.  First fill only $y^q,z^q$ upstairs.
The group downstairs is
$D=C_q(y)\amalg C_q(z)$, with its map to
$C_q$ given by $y\mapsto1$, $z\mapsto r=\zeta^2$.
Its kernel $K$ has Schreier generators
\[
 v_i=y^iz y^{-i-r}\quad(i\in\Z/q),\qquad
           v_0v_rv_{2r}\cdots v_{(q-1)r}=1.              \tag{D.42m}
\]
The displayed relation contains every generator
once and eliminates one; these are the completed
discrete Schreier presentation and maps.
Conjugation by $y$ sends $v_i$ to $v_{i+1}$.
The remaining cusp $x^q=(z^{-1}y^{-1})^q$
rewrites as
\[
             \prod_{j=0}^{q-1}
                    v_{-r-j(r+1)}^{-1}.                \tag{D.42n}
\]
Since $r+1\ne0$, it also uses every index once.
Its abelianization is the negative of the
already imposed sum relation.  Filling it changes
no first homology.  It follows, for the actual
map $K\to\Pi_h$, that
\[
 H_1(\Pi_h,\Q_3)
   =\Q_3[C_q]/\Q_3(1+T+\cdots+T^{q-1})
   \simeq\Q_3[X]/\Phi_q(X).                             \tag{D.42o}
\]
This is the regular augmentation module, with
each irreducible factor occurring once.

Because $q\ne3$, reduction of $\Phi_q$ modulo
three is square-free.  Hensel lifting its
irreducible factors expresses the last algebra
in (D.42o) as a product of finite unramified
extensions of $\Q_3$.
None contains a nonidentity cube root of one:
reduction to the residue field, whose multiplicative
group has order $3^d-1$, first gives residue one;
the same valuation calculation used for $\Z_3$
then applies, since three remains a uniformizer.
The centralizer of $T$ on its cyclic module
(D.42o) is exactly that product algebra.
Thus the commuting operator $b_*$, with
$b_*^3=1$, equals one in every factor.
The torsion-free lattice
$H_1(\Pi_h,\Z_3)$ embeds in (D.42o), so $b_*$
is also the identity on that lattice.
The first part of this proof, applied with $p=3$,
now gives $b=1$ on the full $\Pi_h$.

Finally the quotient filling comparison in A.2
identifies the relative orientation line of
$(N;\langle x^q\rangle,\langle y^q\rangle,
\langle z^q\rangle)$ with the closed orientation
line of $\Pi_h$.
The three cusp maps have the common power
$\lambda$, so $b$ acts on the latter line by
$\lambda$ as well.
Its being the identity forces
$\lambda_q=1$, contradicting the chosen coordinate.
There is no such $q$, proving $\lambda=1$.

For the finite-subgroup assertion, D.32 embeds
every finite subgroup in $S_3\times C_2$.
Its elements therefore have orders $1,2,3$, or $6$.
D.34, D.36, and D.41 give the global sign for
every involution; the preceding proof handles
order three.  An element $v$ of order six is
the product of its commuting powers $v^3$ and
$v^4$, of orders two and three respectively.
Multiplicativity gives the assertion for every
element, with its prescribed $2$-adic sign.

The program \path{code/part_a/pants_order_three.py}
executes the character invariance, both Schreier
index cycles, the cyclic deck matrix and its
cyclotomic polynomial, and the characteristic-three
factorization checks for both primitive cube roots
at every prime $q\equiv1\pmod3$ through $97$.
The arguments above prove the result for every
prime and for the full profinite groups.
\end{proof}
